\begingroup
\raggedright
\GaussContentsEntry{\SourceLeaf{667}\noindent\hangindent=1em Connoissance des tems ou des mouvemens célestes, à l'usage des astronomes et des navigateurs, pour l'an 1810; publiée par le bureau des longitudes. De l'imprimerie impériale 1808. Paris\GaussContentsLeaders{\mbox{.\ .\ .\ .\ .\ .\ .\ .}}\GaussContentsReference{\mbox{1809 June 17}\quad\mbox{Page 515}}}
\GaussContentsEntry{\noindent\hangindent=1em \textsc{von Lindenau}: Results of Investigations into the Diameter of the Sun\GaussContentsLeaders{\mbox{.\ .\ .\ .\ .\ .\ .\ .\ .\ .\ .\ .\ .\ .\ .\ .\ .\ .\ .\ .\ .\ .\ .\ .}}\GaussContentsReference{\mbox{1809 July 17}\quad\mbox{--- 517}}}
\GaussContentsEntry{\noindent\hangindent=1em \textsc{De Rossel}: Voyage de \textsc{Dentrecasteaux}, envoyé à la recherche de la Pérouse. Publié par ordre de Sa Majesté l'Empereur et Roi. Tome second. Paris 1808\GaussContentsLeaders{\mbox{.\ .\ .\ .\ .\ .\ .\ .\ .\ .\ .\ .\ .\ .\ .\ .\ .\ .\ .\ .\ .\ .\ .}}\GaussContentsReference{\mbox{1809 Oct. 21}\quad\mbox{--- 520}}}
\GaussContentsEntry{\noindent\hangindent=1em \textsc{Regner}: Supplementum ad historiam de Parallaxeos Solaris inventione\GaussContentsLeaders{\mbox{.\ .}}\GaussContentsReference{\mbox{1810 Jan. 29}\quad\mbox{--- 523}}}
\GaussContentsEntry{\noindent\hangindent=1em \textsc{von Lindenau}: [Greenwich Observations of the Sun, 1787---1798]\GaussContentsLeaders{\mbox{.\ .\ .}}\GaussContentsReference{\mbox{1810 June 16}\quad\mbox{--- 525}}}
\GaussContentsEntry{\noindent\hangindent=1em Connoissance des tems ou des mouvemens célestes à l'usage des astronomes et des navigateurs, pour l'an 1811, publiée par le bureau des longitudes. Paris 1809\GaussContentsLeaders{\mbox{.\ .\ .\ .\ .\ .\ .\ .\ .\ .\ .\ .\ .\ .\ .\ .\ .\ .\ .\ .\ .}}\GaussContentsReference{\mbox{1810 Sept. 24}\quad\mbox{--- 327}}}
\GaussContentsEntry{\noindent\hangindent=1em \textsc{Knitlmayer}: An Attempt at an Exact Account of the Progression in the Distances of Planets and Satellites from Their Central Bodies. Brünn 1808\GaussContentsLeaders{\mbox{.\ .\ .\ .\ .\ .\ .\ .\ .\ .\ .\ .\ .\ .\ .\ .\ .\ .\ .\ .\ .\ .}}\GaussContentsReference{\mbox{1810 Nov. 10}\quad\mbox{--- 528}}}
\GaussContentsEntry{\noindent\hangindent=1em \textsc{Bessel}: Investigations into the Apparent and True Orbit of the Great Comet That Appeared in 1807. Königsberg\GaussContentsLeaders{\mbox{.\ .\ .\ .\ .\ .\ .\ .\ .}}\GaussContentsReference{\mbox{1811 Feb. 16}\quad\mbox{--- 530}}}
\GaussContentsEntry{\noindent\hangindent=1em \textsc{Schulze}: Minister at Polenz: An Account of the System of the World. A Guide for Teaching Astronomy in Schools. Leipzig 1811\GaussContentsLeaders{\mbox{.\ .}}\GaussContentsReference{\mbox{1811 Feb. 23}\quad\mbox{--- 533}}}
\GaussContentsEntry{\noindent\hangindent=1em \textsc{Seyffer}: Super longitudine geographica speculae astronomicae regiae, quae Monachii est, ex triginta septem defectionibus Solis observatis et ad calculos revocatis nunc primum definita. Munich 1810\GaussContentsLeaders{\mbox{.\ .\ .\ .\ .}}\GaussContentsReference{\mbox{1811 May 18}\quad\mbox{--- 533}}}
\GaussContentsEntry{\noindent\hangindent=1em \textsc{Nouet}: Description de l'Egypte. Tome premier à Paris de l'imprimerie impériale 1809. Mémoire premier. Observations astronomiques, faites en Egypte pendant les années VI, VII et VIII\GaussContentsLeaders{\mbox{.\ .\ .\ .\ .\ .\ .\ .\ .}}\GaussContentsReference{\mbox{1811 May 18}\quad\mbox{--- 534}}}
\GaussContentsEntry{\noindent\hangindent=1em \textsc{Pasquich}: Epitome elementorum astronomiae sphaerico-calculatoriae.\par\noindent\hangindent=1em\hspace{1em}Pars prima: Elementa theoretica astronomiae sphaerico-calculatoriae.\par\noindent\hangindent=1em\hspace{1em}Pars secunda: Elementa practica astronomiae sphaerico-calculatoriae. Vienna 1811\GaussContentsLeaders{\mbox{.\ .\ .\ .\ .\ .\ .\ .\ .\ .\ .\ .\ .\ .\ .\ .\ .\ .\ .\ .\ .\ .\ .}}\GaussContentsReference{\mbox{1811 Sept. 9}\quad\mbox{--- 535}}}
\GaussContentsEntry{\noindent\hangindent=1em \textsc{Klüber}: The Mannheim Observatory. Mannheim, Heidelberg 1811\GaussContentsLeaders{\mbox{.\ .\ .}}\GaussContentsReference{\mbox{1811 Dec. 2}\quad\mbox{--- 536}}}
\GaussContentsEntry{\noindent\hangindent=1em Connoissance des tems, etc.\ .\ .\ . pour l'an 1812. Paris 1810\GaussContentsLeaders{\mbox{.\ .\ .\ .\ .\ .}}\GaussContentsReference{\mbox{1811 Dec. 14}\quad\mbox{--- 537}}}
\GaussContentsEntry{\noindent\hangindent=1em D'Aquila: Decouverte de l'orbite de la terre du point central de l'orbite du soleil, leur situation et leur forme; de la section du zodiaque, par le plan de l'équateur, et du mouvement concordant des deux globes. Paris 1806\GaussContentsLeaders{\mbox{.\ .\ .\ .\ .\ .\ .\ .\ .\ .\ .\ .\ .\ .\ .\ .\ .\ .\ .\ .\ .\ .}}\GaussContentsReference{\mbox{1812 June 27}\quad\mbox{--- 541}}}
\GaussContentsEntry{\noindent\hangindent=1em \textsc{Harding}: Star Atlas. Instalments III and IV.\GaussContentsLeaders{\mbox{.\ .\ .\ .\ .\ .\ .\ .\ .\ .}}\GaussContentsReference{\mbox{1812 Oct. 3}\quad\mbox{--- 541}}}
\GaussContentsEntry{\noindent\hangindent=1em \textsc{Gerling}: Methodi projectionis orthographicae usum ad calculos parallacticos facilitandos explicavit, simulque eclipsin solarem die VII. Sept. 1820 apparituram hoc modo tractatam mappaque geographica illustratam tamquam exemplum proposuit. Göttingen\GaussContentsLeaders{\mbox{.\ .\ .\ .\ .\ .\ .\ .\ .\ .\ .}}\GaussContentsReference{\mbox{1812 Nov. 28}\quad\mbox{--- 542}}}
\GaussContentsEntry{\noindent\hangindent=1em \textsc{Carlini}: Esposizione di un nuovo methodo di costruire le tavole astronomiche applicato alle tavole del sole. Milan 1810.\GaussContentsLeaders{\mbox{.\ .\ .\ .\ .\ .\ .}}\GaussContentsReference{\mbox{1812 Dec. 3}\quad\mbox{--- 543}}}
\GaussContentsEntry{\noindent\hangindent=1em \textsc{De Zach}: Nouvelles tables d'aberration et de nutation pour quatorze cent \SourceLeaf{668}quatre étoiles, avec une table générale d'aberration pour les planètes et les comètes, précédées d'une instruction qui renferme l'explication de l'usage de ces tables, suivies de plusieurs nouvelles tables destinées à faciliter les calculs astronomiques. Marseille\GaussContentsLeaders{\mbox{.\ .\ .\ .\ .\ .\ .\ .\ .}}\GaussContentsReference{\mbox{1812 Dec. 7}\quad\mbox{Page 546}}}
\GaussContentsEntry{\noindent\hangindent=1em \textsc{Bode}: Astronomical Yearbook for the Year 1814. Berlin 1811\par\noindent\hangindent=1em\hspace{1em}Astronomical Yearbook for the Year 1815. Berlin 1812\GaussContentsLeaders{\mbox{.\ .\ .\ .\ .}}\GaussContentsReference{\mbox{1813 Feb. 6}\quad\mbox{--- 548}}}
\GaussContentsEntry{\noindent\hangindent=1em Tables astronomiques publiées par le bureau des longitudes de France.\par\noindent\hangindent=1em \textsc{Burckhardt}: Tables de la lune. Paris 1812.\GaussContentsLeaders{\mbox{.\ .\ .\ .\ .\ .\ .\ .\ .\ .\ .}}\GaussContentsReference{\mbox{1813 March 8}\quad\mbox{--- 553}}}
\GaussContentsEntry{\noindent\hangindent=1em Connaissance des tems etc.\ .\ .\ . pour l'an 1813. Paris 1811.\par\noindent\hangindent=1em Connaissance des tems etc.\ .\ .\ . pour l'an 1814. Paris 1812\GaussContentsLeaders{\mbox{.\ .\ .\ .\ .\ .}}\GaussContentsReference{\mbox{1813 June 3}\quad\mbox{--- 554}}}
\GaussContentsEntry{\noindent\hangindent=1em \textsc{Carlini} e \textsc{Brioschi}: Effemeridi astronomische di Milano per l'anno bisestile 1812. Milano 1811.\par\noindent\hangindent=1em\hspace{1em}Effemeridi astronomiche di Milano per l'anno 1813\GaussContentsLeaders{\mbox{.\ .\ .\ .\ .\ .\ .\ .}}\GaussContentsReference{\mbox{1813 June 14}\quad\mbox{--- 556}}}
\GaussContentsEntry{\noindent\hangindent=1em \textsc{De Lindenau}: Investigatio nova orbitae a Mercurio circa Solem descriptae: accedunt tabulae Planetae ex elementis recens correctis et ex theoria gravitatis ill. de la Place constructae\GaussContentsLeaders{\mbox{.\ .\ .\ .\ .\ .\ .\ .\ .\ .\ .\ .}}\GaussContentsReference{\mbox{1813 July 5}\quad\mbox{--- 559}}}
\GaussContentsEntry{\noindent\hangindent=1em \textsc{Bessel}: Determination of Precession [Berlin Prize Essay]\GaussContentsLeaders{\mbox{.\ .\ .\ .\ .\ .}}\GaussContentsReference{\mbox{1814 Jan. 31}\quad\mbox{--- 561}}}
\GaussContentsEntry{\noindent\hangindent=1em \textsc{Pond}: Observation of the Summer Solstices of 1812 and 1813 and the Winter Solstice of 1812 at the Royal Observatory at Greenwich. Philosophical Transactions of the Royal Society of London for the Jear 1813\GaussContentsLeaders{\mbox{.\ .\ .\ .\ .\ .\ .\ .\ .\ .\ .\ .\ .\ .\ .\ .\ .\ .\ .\ .\ .\ .\ .}}\GaussContentsReference{\mbox{1814 Feb. 14}\quad\mbox{--- 561}}}
\GaussContentsEntry{\noindent\hangindent=1em \textsc{Bode}: Astronomical Yearbook for the Year 1816\GaussContentsLeaders{\mbox{.\ .\ .\ .\ .\ .\ .\ .\ .\ .}}\GaussContentsReference{\mbox{1814 March 10}\quad\mbox{--- 563}}}
\GaussContentsEntry{\noindent\hangindent=1em \textsc{Carlini}: Effemeridi astronomiche di Milano per l'anno 1814\GaussContentsLeaders{\mbox{.\ .\ .\ .\ .\ .}}\GaussContentsReference{\mbox{1814 March 24}\quad\mbox{--- 566}}}
\GaussContentsEntry{\noindent\hangindent=1em \textsc{Mollweide}: Commentationes mathematico-philologicae tres, sistentes explicationem duorum locorum difficilium, alterius Virgilii, alterius Platonis itemque examinationem duorum mensurarum praeceptorum Columellae. Leipzig 1813'\GaussContentsLeaders{\mbox{.\ .\ .\ .\ .\ .\ .\ .\ .\ .\ .\ .\ .\ .\ .\ .\ .\ .}}\GaussContentsReference{\mbox{1814 May 2}\quad\mbox{--- 567}}}
\GaussContentsEntry{\noindent\hangindent=1em \textsc{De Zach}: L'attraction des montagnes, et ses effets sur les fils à plomb ou sur les niveaux des instrumens d'astronomie,\ .\ .\ .\ . Avignon 1814\GaussContentsLeaders{\mbox{.\ .}}\GaussContentsReference{\mbox{1814 Dec. 31}\quad\mbox{--- 569}}}
\GaussContentsEntry{\noindent\hangindent=1em \textsc{Carlini}: Effemeridi astronomiche di Milano per l'anno 1815\GaussContentsLeaders{\mbox{.\ .\ .\ .\ .\ .}}\GaussContentsReference{\mbox{1815 Jan. 19}\quad\mbox{--- 574}}}
\GaussContentsEntry{\noindent\hangindent=1em \textsc{Bode}: Astronomical Yearbook for the Year 1817\GaussContentsLeaders{\mbox{.\ .\ .\ .\ .\ .\ .\ .\ .\ .}}\GaussContentsReference{\mbox{1815 Feb. 18}\quad\mbox{--- 576}}}
\GaussContentsEntry{\noindent\hangindent=1em Connaissance des tems etc.\ .\ .\ . pour l'an 1816\GaussContentsLeaders{\mbox{.\ .\ .\ .\ .\ .\ .\ .\ .\ .\ .}}\GaussContentsReference{\mbox{1815 March 11}\quad\mbox{--- 581}}}
\GaussContentsEntry{\noindent\hangindent=1em \textsc{Soldner}: A New Method for Reducing Observed Azimuths\GaussContentsLeaders{\mbox{.\ .\ .\ .\ .}}\GaussContentsReference{\mbox{1815 March 23}\quad\mbox{--- 587}}}
\GaussContentsEntry{\noindent\hangindent=1em \textsc{Stefenelli}: Contributions to the Calculation of Observed Azimuths\GaussContentsLeaders{\mbox{.\ .\ .\ .\ .\ .}}\GaussContentsReference{\mbox{1815 March 23}\quad\mbox{--- 587}}}
\GaussContentsEntry{\noindent\hangindent=1em \textsc{Piazzi}: Praecipuarum stellarum inerrantium positiones etc. Palermo 1814\GaussContentsLeaders{\mbox{.}}\GaussContentsReference{\mbox{1815 June 10}\quad\mbox{--- 590}}}
\GaussContentsEntry{\noindent\hangindent=1em \textsc{Müller}: Calculation of Solar Eclipses\GaussContentsLeaders{\mbox{.\ .\ .\ .\ .\ .\ .\ .\ .\ .\ .\ .}}\GaussContentsReference{\mbox{1816 Feb. 3}\quad\mbox{--- 591}}}
\GaussContentsEntry{\noindent\hangindent=1em \textsc{Carlini}: Effemeridi astronomiche di Milano per l'anno bisestile 1816\GaussContentsLeaders{\mbox{.\ .\ .}}\GaussContentsReference{\mbox{1816 March 23}\quad\mbox{--- 592}}}
\GaussContentsEntry{\noindent\hangindent=1em \textsc{Harding}: Star Atlas. Instalment V.\GaussContentsLeaders{\mbox{.\ .\ .\ .\ .\ .\ .\ .\ .\ .\ .\ .\ .\ .}}\GaussContentsReference{\mbox{1817 Jan. 9}\quad\mbox{--- 593}}}
\GaussContentsEntry{\noindent\hangindent=1em Connaissance des tems etc.\ .\ .\ . pour l'an 1818\GaussContentsLeaders{\mbox{.\ .\ .\ .\ .\ .\ .\ .\ .\ .\ .}}\GaussContentsReference{\mbox{1817 April 5}\quad\mbox{--- 595}}}
\GaussContentsEntry{\noindent\hangindent=1em \textsc{Schumacher}: De latitudine speculae Manhemiensis. Copenhagen 1816\GaussContentsLeaders{\mbox{.\ .}}\GaussContentsReference{\mbox{1817 May 12}\quad\mbox{--- 597}}}
\GaussContentsEntry{\noindent\hangindent=1em Tavole delle parallassi, di altezza e di latitudine calcolate dagli astronomi dell' osservatorio\ .\ .\ . nel collegio Romano. Rome 1816\GaussContentsLeaders{\mbox{.\ .\ .\ .\ .\ .\ .}}\GaussContentsReference{\mbox{1817 July 7}\quad\mbox{--- 599}}}
\GaussContentsEntry{\noindent\hangindent=1em \textsc{Carlini}: Effemeridi astronomiche di Milano per l'anno 1817. Milan 1816\GaussContentsLeaders{\mbox{.}}\GaussContentsReference{\mbox{1817 Sept. 8}\quad\mbox{---\textperiodcentered{} 601}}}
\endgroup
